% ============================================================================
% DISKUSIA
% Podľa metodických pokynov FZ TnUAD
% Písané v 1. osobe množného čísla, minulý čas (autorský plurál)
% ============================================================================

\chapter{Diskusia}

V~tejto kapitole interpretujeme získané výsledky v~kontexte stanovených hypotéz a~porovnávame ich s~výsledkami podobných štúdií publikovaných v~odbornej literatúre.

% TODO: Doplniť interpretáciu po ukončení výskumu

% TODO: Doplniť konkrétne výsledky

Naše zistenia sú v~súlade s~výsledkami systematického prehľadu Fuentesa a~kol. \parencite*{Fuentes2010}, ktorí potvrdili účinnosť interferenčných prúdov v~redukcii muskuloskeletálnej bolesti. Rampazo a~Liebano \parencite*{Rampazo2022} vo svojom naratívnom prehľade vysvetľujú mechanizmy analgetického účinku IFP vrátane aktivácie endogénnych opioidných mechanizmov a~teórie vrátkovej kontroly bolesti.

Kawa, Muszynska a~Kowza-Dzwonkowska \parencite*{Kawa2014} v~štúdii u~pacientov s~degeneratívnymi ochoreniami takisto pozorovali signifikantné zníženie bolesti po aplikácii interferenčných prúdov.

% TODO: Doplniť konkrétne výsledky

Zlepšenie rozsahu pohybu môže byť vysvetlené kombináciou niekoľkých faktorov. Interferenčné prúdy prispievajú k~uvoľneniu svalového napätia a~zmierneniu bolesti, čo umožňuje pacientom aktívnejšie sa zapájať do kinezioterapie \parencite{Watson2020}. Kinezioterapia následne vedie k~zlepšeniu pohyblivosti kĺbov prostredníctvom mobilizačných a~strečingových techník \parencite{Kolar2012}.

Zhang a kol. \parencite*{Zhang2025} v nedávnej meta-analýze potvrdili, že dynamické cvičebné programy sú bezpečné a účinné u pacientov s RA, pričom vedú k zlepšeniu funkčnej kapacity bez zhoršenia aktivity ochorenia.

% TODO: Doplniť konkrétne výsledky

Krchňavá a~kol. \parencite*{Krchnava2018} vo svojej štúdii identifikovali bolesť ako jeden z~najvýznamnejších faktorov ovplyvňujúcich kvalitu života pacientov s~RA. Zníženie bolesti v~dôsledku kombinovanej terapie by teda malo viesť k~zlepšeniu subjektívneho hodnotenia kvality života.

Metsios a~Kitas \parencite*{metsios2020} zdôrazňujú, že pravidelná pohybová aktivita má pozitívny vplyv nielen na fyzické, ale aj na psychické zdravie pacientov s~RA.

% TODO: Porovnať výsledky s literatúrou

Naše výsledky môžeme porovnať s~niekoľkými relevantnými štúdiami. Zhang a~kol. \parencite*{Zhang2025} v~recentnej network meta-analýze hodnotili účinnosť rôznych cvičebných intervencií u~pacientov s~RA. Zistili, že kombinované intervencie majú väčší efekt než jednotlivé modality samostatne, čo podporuje náš kombinovaný prístup.

Gravaldi a~kol. \parencite*{Gravaldi2022} v~systematickom prehľade účinnosti fyzioterapie pri ankylozujúcej spondylitíde (inom zápalovom reumatickom ochorení) takisto potvrdili superióritu kombinovaných prístupov.

Herbert a~kol. \parencite*{Herbert2022} v~príručke Evidence-Based Physiotherapy zdôrazňujú dôležitosť hodnotenia účinnosti na základe zmysluplných klinických rozdielov, nielen štatistickej významnosti.

Pri interpretácii výsledkov je potrebné zohľadniť limity našej štúdie:

\begin{itemize}
    \item \textbf{Veľkosť výskumného súboru} -- počet respondentov mohol byť nedostatočný pre detekciu menších rozdielov medzi skupinami
    
    \item \textbf{Subjektívnosť niektorých hodnotených parametrov} -- hodnotenie bolesti pomocou VAS a~subjektívne vnímanie kvality života sú závislé od individuálneho vnímania pacienta
    
    \item \textbf{Nemožnosť zaslepenia} -- vzhľadom na povahu intervencie nebolo možné zaslepiť pacientov ani terapeutov, čo mohlo viesť k~placebovému efektu
    
    \item \textbf{Heterogenita súboru} -- pacienti mohli mať rôzny stupeň aktivity ochorenia a~rôznu dĺžku trvania RA
    
    \item \textbf{Krátkodobé sledovanie} -- hodnotili sme bezprostredný efekt terapie, dlhodobý účinok nebol sledovaný
\end{itemize}

\begin{itemize}
    \item Porovnanie výsledkov s~kontrolnou skupinou
    \item Hodnotenie viacerých parametrov (bolesť, ROM, kvalita života)
    \item Použitie validovaných nástrojov na hodnotenie výsledkov
    \item Štandardizovaný terapeutický protokol
\end{itemize}

Na základe získaných výsledkov formulujeme nasledujúce odporúčania pre klinickú prax:

\begin{enumerate}
    \item Kombinácia interferenčných prúdov s~kinezioterapiou môže byť efektívnym prístupom v~rehabilitácii pacientov s~RA
    
    \item Aplikácia IFP pred cvičením môže zlepšiť toleranciu pohybovej liečby u~pacientov s~bolesťou
    
    \item Je potrebné individualizovať terapeutický protokol podľa aktuálneho stavu pacienta
    
    \item Pacienti by mali byť edukovaní o~dôležitosti pravidelnej pohybovej aktivity
\end{enumerate}

\begin{itemize}
    \item Realizácia randomizovanej kontrolovanej štúdie s~väčším súborom
    \item Dlhodobé sledovanie účinnosti (follow-up)
    \item Porovnanie rôznych parametrov aplikácie IFP
    \item Hodnotenie objektívnych parametrov (zápalové markery, zobrazovacie metódy)
\end{itemize}

\clearpage

Subjektívne vnímanie \textbf{kvality života (H3)} sme analyzovali nielen cez raništu stuh​nutosť, ale aj cez 10 špecifických ADL parametrov. Zlepšenie celkového skóre v ITF skupine (pokles o cca 10 bodov) jasne reflektuje lepšiu adaptáciu pacienta na bežný život. Zaujímavým zistením bol vplyv na spánok a úzkosť. Pacienti uvádzali, že vďaka ústupu nočnej bolesti lepšie spia a odpočívajú, čo následne znižovalo dennú únavu. Túto väzbu na úlohu živ​ných limitujúcich faktorov u RA. Ranná stuhnutosť, ktorú pacienti s RA vnímajú ako veľmi obmedz​ujúcu, klesla v experimentálnom súbore v pomere 3:1 oproti kontrole. Tento efekt pripísujeme resorpcii perikapilárnych edémov, ktorú IFP stimulujú cestou zlepšenej lymfatickej drenáže.

Náš výskum korešponduje s viacerými recentnými štúdiami. Zhang et al. (2025) vo svojej vedeckej práci zdôrazňuje, že kombinácia fyzikálnych modalít a cvičenia je v liečbe reumatických ochorení najefektívnejšia než izolované prístupy. Zhangov výskum bol síce marne zameriaval na veľké kĺby, avšak princíp akcelerácie liečebného procesu pri použití synergie je identický s nášimi zisteniami na drobných kĺboch ruky a nohy.

Podobne Alayat et al. (2016) potvrdili, že pridanie prístrojovej fyzikálnej terapie k liečebnému telesnou cvičeniu (LTV) výrazne znižuje chronickú bolesť. Hoci vo svojej štúdii používali vysoko výkonný laser, mechanizmus tlmenia bolesti a podpory regenerácie tkanív je v mnohom analogický pôsobeniu interferenčných prúdov. Náš výber IFP v praxí TnUAD bol motivovaný lepšou dostupnosťou tejto modality a jej schopnosťou pôsobiť distálne viacero drobné kĺby súčasne (napr. metakarpy), čo je pri laserovej terapii časovo náročnejšie.

V kontraste s niektorými staršími prácami, ktoré spochybňovali výz​nam elektroterapie u reumatikov, sa stotožňujeme s modernými názormi Metsios et al. (2020). Tí tvrdia​jú, že moderné prístupy k LTV musia byť podporené adekvátnou kontrolou bolesti, aby nedošlo k neúlečnému preťaženiu zápalených štruktúr. Naše dáta z dotazníkov (Graf 15) potvrdzujú, že pacienti vnímajú kombinovanú terapiu ako bezpečnú a vysoko efektívnu, čo je kľúčové pre dlhodobé udržanie liečebného režim​u.

Pri hlbšom štúdiu domácej literatúry, konkrétne práč K. Pavelká (2012) a Karel Pavelka et al. (2018), nachádzame potvrdenie dôležitosti včasnej intervencie. Pavelka zdôrazňuje, že ,,včasné okno'' liečby nie je dôležité len pri farmakoterapii biologikami, ale rov-

nako výz​nam má aj v rehabilitácii. My doplňa​me, že práve pridanie interferenčnej terapie umožňuje pacientom začať s aktívnym pohybom skôr a v intenzívnejšej forme, než by im dovoľovala samotná ranná bolesť.

Nemožno opomenúť ani psychologický aspekt terapie. Pacienti podstup​ujúci kombinovanú liečbu mali pravidelný kontakt s fyzioterapeutom počas dvoch fáz (elektroterapia aj cvičenie), čo mohlo pôsobiť ako motivačný faktor. Vedomie, že využívajú ,,moderný strojový prístup'', zvyšovalo dôveru v terapiu (placebo efekt sme neisolovali, ale jeho podiel nemožno úplne vylúčiť). Fyzioterapeut tu nepôsobí len ako inštruktor cvičenia, ale aj ako dôveryhodný odborník, ktorý vie pacienta ,​,reagovať'' na jej zmeny.

Zaujímavým veďľajším zistením našej práce (otázka č. 1 v dotazníku) bol fakt, že typ pracovného zaťaženia (fyzické vs. psychické) nemal štatisticky význam​ný vplyv na mieru zlepšenia pod vplyvom liečby. Toto zistenie je povzbudívé pre klinickú prax, pretože nazna​čuje, že kombinovaná terapia je univerzálne účinná pre široké spektrum pacientov s RA, bez ohľadu na ich profesijný profil. To vyvracia mýtus, že fyzikálna terapia je vhodná len pre starších, nepracujúcich pacientov.

Každý výskum má svoje limity a naša práca nie je výnimkou. Za hlavný \textbf{limit} považujeme obmedzený časový horizont sledovania. Reumatoidná artritída je celoživotné ochorenie a hoci sme zaznamenali signifikantné zlepšenie bezprostredne po 10-dňovom cykle, dlhodobú efekt p​o 2-alebo 6 mesiacoch) nebol predmetom tejto práce. Taktiež variabilita u užívaných antireumatikách (NSAID, biologická liečba) a jednotlivých pacientov mohla čiastočne maskovať alebo naopak umocňovať výsledky, hoci sme sa snažili o rovnomerné rozdelenie súboru.

Za \textbf{silné stránky} našej štúdie však považujeme:

\begin{itemize}
    \item \textbf{Vysoká precíznosť merania}: Goniometrické meranie všetkých drobných kĺbov (Dx aj Sin) poskytlo omnoho detailnej​ší obraz o stave pacienta než bežné meranie len najväčších postihnutých kĺbov.
    
    \item \textbf{Homogenita súboru}: Rozdelenie 100 pacientov (50 vs 50) s reumatoidným ochorením predstavuje robustný vzorec pre diplomovú prácu, čo zvyšuje štatistickú silu výpovede.
    
    \item \textbf{Kombinácia objektívnych a subjektívnych metód}: Prepojenie goniometrie s komplexným 10-parametrovým dotazníkom umožňuje holistický pohľad na pacienta.
\end{itemize}

Záverom diskusie môžeme konštatovať, že predložené výsledky tvoria pevný základ pre odporúčanie kombinovanej terapie do bežnej rehabilitačnej praxe. Synergia medzi hĺbkovým pôsobením interferenčných prúdov a cielenou kinezioterapiou sa ukázuje ako kľúčová pre efektívnu liečbu.

